% !PRINTED 175
\weberpage{175}
%<<CONT
variante $J$ nach dem Satze 4. als rationale Function von $\Theta$
darstellen, so kann sich diese nicht ändern, wenn eine der Substitutionen
$(\Theta, \Theta_1), (\Theta, \Theta_2), \dots$ ausgeführt wird, und diese Function
ist also rational durch die Co\"efficienten von $\Phi(t)$ ausdrückbar.
Ob diese Darstellung freilich durch ganze Functionen
möglich ist, würde bei diesem Beweise unentschieden bleiben.

Unter den Functionen $\omega$ der Variablen $x$ sind auch die
Variablen $x$ selbst enthalten, und die Frage, von der wir ausgegangen
sind, ob die Variablen $x$ als algebraische Functionen
der Invarianten angesehen werden können, ist damit bejahend
entschieden.

Die Aufgabe, die Variablen $x$ als algebraische Functionen
der Invarianten der Gruppe $S$ darzustellen, also die Bestimmung
des Körpers $\Omega(\Theta)$ heisst nach \gesperrt{F. Klein} das \gesperrt{Formenproblem
der Gruppe $S$}\footnote{Vorlesungen über das Ikosaëder, S.\,123. (Leipzig 1884.)}. Ist die Anzahl der Variablen $n$, so nennen
wir das Formenproblem von der $n^{\text{ten}}$ Dimension.

Der Körper $\Omega(\Theta)$ ist ein durch die Gruppe $S$ völlig bestimmter
algebraischer Körper über $\Omega$. Er ist ein Normalkörper,
denn nach 4. sind die conjugirten Grössen $\Theta, \Theta_1, \Theta_2, \dots \Theta_{\mu-1}$
alle im Körper $\Omega(\Theta)$ selbst enthalten. Die Gleichung $\Phi(t) = 0$
ist eine Normalgleichung und ist die \gesperrt{Galois'sche Resolvente
des Formenproblems} (Bd.\ I, \S.\,145).

\gesperrt{5. Die Galois'sche Gruppe des Formenproblems,
d. h. die Galois'sche Gruppe der Gleichung
$\Phi(t) = 0$, ist mit der Gruppe $S$ isomorph.}

Um dies nachzuweisen, bezeichnen wir mit $A, B$ zwei Substitutionen
aus $S$ und mit $AB = C$ die daraus zusammengesetzte
Substitution. Ist nun
\[ \Theta_1 = \Theta[A(x)],\ \Theta_2 = \Theta[B(x)],\ \Theta_3 = \Theta[C(x)], \]
so ist die Substitution
\[ (\Theta, \Theta_2) = (\Theta_1, \Theta_3), \]
und folglich
\[ (\Theta, \Theta_1)(\Theta, \Theta_2) = (\Theta, \Theta_1)(\Theta_1, \Theta_3) = (\Theta, \Theta_3), \]
d. h. die Gruppe der Substitutionen $(\Theta, \Theta_1), (\Theta, \Theta_2), \dots$ ist mit
der Gruppe der $A, B, \dots$ isomorph.

Nehmen wir statt der Function $\Theta$ eine Function $\eta$, die nicht
lauter verschiedene Werthe hat, sondern die Substitutionen eines
%CONT>>
